\metatitle{Z-polynomials of Schubert matroids}

\metadescription{Combinatorial formulas for the Z-polynomials of Schubert and
lattice path matroids: Motzkin paths with rook placements, words, tableaux,
pattern-avoiding permutations, ribbon matroids, a row recursion and the
real-rootedness question.}
\metakeywords{Z-polynomial,matroid Kazhdan--Lusztig polynomials,Schubert
matroids,lattice path matroids,Motzkin paths,rook placements,gamma-positivity,
real-rooted polynomials}


\section[schubertMatroidZ]{Z-polynomials of Schubert matroids}

Let $M$ be a loopless matroid of rank $n$, and let $P_M(t)$ be its
\hyperref[matroidKazhdanLusztigPolynomial]{Kazhdan--Lusztig polynomial}.
The \defin{$Z$-polynomial}\label{matroidZPolynomial} of $M$, introduced by
\name{Nicholas Proudfoot}, \name{Yuan Xu} and \name{Ben Young}
\cite{ProudfootXuYoung2018}, is
\[
  Z_M(t) \coloneqq \sum_{F \text{ a flat}} t^{\operatorname{rk} F}\, P_{M/F}(t).
\]
It is palindromic of degree $n$, so there are unique numbers $\gamma_k$ with
\[
  Z_M(t)=\sum_{k\geq0}\gamma_k\, t^k(1+t)^{n-2k},
\]
and we write $\Gamma_M(y)\coloneqq\sum_k\gamma_k y^k$ for the
\hyperref[gammaPolynomial]{gamma polynomial}.  The $\gamma_k$ are nonnegative
for every matroid \cite[Thm.~1.8]{FerroniMatherneStevensVecchi2022x}.
The polynomial $Z_M$ depends only on the lattice of flats, so parallel
elements may be removed, and it is multiplicative on direct sums; in particular
each coloop contributes a factor $1+t$.  For the circuit $U_{n,n+1}$,
the polynomial $Z_{U_{n,n+1}}(t)$ is the \hyperref[narayanaNumbers]{Narayana polynomial}
\cite{ProudfootXuYoung2018}
\[
  N_{n+1}(t)\coloneqq
  \sum_{k=0}^{n}\frac{1}{n+1}\binom{n+1}{k}\binom{n+1}{k+1}t^k
  =\sum_{k\geq0}\binom{n}{2k}\mathrm{Cat}_k\, t^k(1+t)^{n-2k}.
\]

Real-rootedness of $Z_M(t)$ was conjectured for all matroids
\cite{ProudfootXuYoung2018}.  It holds for sparse paving matroids
\cite{XieZhang2026SparsePaving} but fails for the extended binary Golay matroid
\cite{Eberhardt2026x}.  For Schubert matroids the question is open; see the
\hyperref[conjNestedMatroidZ]{nested-matroid conjecture}.  This page collects
combinatorial descriptions of $Z_M$ and $\Gamma_M$ for Schubert matroids and,
more generally, lattice path matroids.  Unless stated otherwise, the results
are from unpublished work of \name{Per Alexandersson}
\cite{Alexandersson2026SchubertMatroidZ}.


\subsection[schubertMatroidZBoard]{Schubert matroids and the board of squares}

Let $\lambda=(\lambda_1\geq\dotsb\geq\lambda_n\geq1)$ be a partition with $n$
positive parts.  The Schubert matroid $M_\lambda$ on $[n+\lambda_1]$ has as
bases the sets $\{j_1<\dotsb<j_n\}$ with $j_a\geq a+\lambda_1-\lambda_a$ for
all $a$.  This is the \hyperref[schubertMatroidDefinition]{Schubert matroid}
of $I=\{a+\lambda_1-\lambda_a\}$, or equivalently the
\hyperref[latticePathMatroids]{lattice path matroid} of the Ferrers shape
$\lambda$.  It has rank $n$ and no loops or coloops.  Zero parts would add
coloops, each multiplying $Z$ by $1+t$.  The elements
$1,\dotsc,1+\lambda_1-\lambda_2$ are parallel, so $Z_{M_\lambda}$ does not
depend on $\lambda_1$.

The \defin{board of squares}\label{schubertMatroidZBoardOfSquares} $B$ of
$\lambda$ consists of the boxes of $\lambda$ that are the northwest corner of a
$2\times2$ square inside $\lambda$.  Its row $j$, for $1\leq j<n$, has
$w_j=\lambda_{j+1}-1$ boxes.  A \defin{placement} on $B$ is a set of boxes
$(j_1,c_1),\dotsc,(j_s,c_s)$ of $B$ with $j_1<\dotsb<j_s$ and
$c_1>\dotsb>c_s$, that is, a
\hyperref[nonNestingRookPlacement]{non-nesting rook placement} on $B$.  Give a
rook in column $c$ the weight $x_c$, and for $E\subseteq[n-1]$ let
\[
  W_B(E)\coloneqq\sum_{\text{placements in rows } \subseteq E}\ \prod_{i} x_{c_i},
\]
including the empty placement of weight one.  The identities below hold with
these column weights, and setting all $x_c=1$ gives $Z_{M_\lambda}(t)$ and
$\Gamma_{M_\lambda}(y)$.

Placements are the same as partitions inside $B$.  Given a placement, let
$\eta_j$ be the column of the first rook in a row $i\geq j$, or zero if there
is none, and put $\eta_n=0$.  Then $\eta$ is a partition contained in $B$,
and the rooks sit exactly at its corners, the boxes $(j,\eta_j)$ with
$\eta_j>\eta_{j+1}$.  Conversely, every partition $\eta\subseteq B$ arises
from exactly one placement.  Thus
\[
  W_B(E)=\sum_{\substack{\eta\subseteq B\\ \text{corners in rows of } E}}
  \ \prod_{\text{corners } (j,\eta_j)} x_{\eta_j}.
\]
The partitions $\eta\subseteq B$ are the lattice paths inside the Ferrers shape
$B$, traced from the bottom left to the top right, and a corner of $\eta$ is
where this path turns from east to north.  Through the
correspondences on the page on \hyperref[latticePathMatroids]{lattice path
matroids}, they are also the bases of the lattice path matroid of $B$ and the
linear extensions of its width-two poset.

\begin{example}
For $\lambda=(3,3,2,1)$ the board has rows of lengths $2,1,0$.  Below, the
board is shaded, the partition $\eta$ is blue and the rooks are dots.  The
five placements are the empty one, $(1,1)$, $(1,2)$, $(2,1)$ and
$\{(1,2),(2,1)\}$.
\begin{figure}
\begin{ytableau}
*(lightgray) & *(lightgray) & \\
*(lightgray) & & \\
 & \\
 \\
\end{ytableau}
\begin{ytableau}
*(lightblue) \bullet & *(lightgray) & \\
*(lightgray) & & \\
 & \\
 \\
\end{ytableau}
\begin{ytableau}
*(lightblue) & *(lightblue) \bullet & \\
*(lightgray) & & \\
 & \\
 \\
\end{ytableau}
\begin{ytableau}
*(lightblue) & *(lightgray) & \\
*(lightblue) \bullet & & \\
 & \\
 \\
\end{ytableau}
\begin{ytableau}
*(lightblue) & *(lightblue) \bullet & \\
*(lightblue) \bullet & & \\
 & \\
 \\
\end{ytableau}
\end{figure}
The corresponding partitions are $\eta=000$, $100$, $200$, $110$ and $210$.
For instance $W_B(\{1,2\})=1+2x_1+x_2+x_1x_2$, while $W_B(\{2\})=1+x_1$.
\end{example}


\subsection[schubertMatroidZMotzkin]{Motzkin paths and rook placements}

A \hyperref[motzkinPolynomialsRealRooted]{Motzkin path} of length $n$ starts
and ends at height zero, takes steps $U,H,D$ of increments $1,0,-1$ and stays
weakly above zero.  Let $h_j(P)$ be the height after step $j$ and $k(P)$ the
number of up-steps.

\begin{theorem}[Alexandersson, 2026]\label{schubertMatroidZRookMotzkin}
For a Motzkin path $P$ of length $n$, put
$E(P)=\{j\in[n-1]: h_j(P)>0\}$.  Then
\[
  \Gamma_{M_\lambda}(y)=\sum_{P} y^{k(P)}\,W_B(E(P)),
  \qquad
  Z_{M_\lambda}(t)=\sum_{P} t^{k(P)}(1+t)^{n-2k(P)}\,W_B(E(P)).
\]
Thus $\gamma_k$ counts pairs of a Motzkin path with $k$ up-steps and a placement
whose rooks all lie in rows where the path has positive height.  To count
$[t^i]Z_{M_\lambda}$, additionally mark $i-k(P)$ of the horizontal steps.
\end{theorem}

In terms of the partitions $\eta\subseteq B$, the coefficient $\gamma_k$
counts the pairs of a partition $\eta\subseteq B$ and a Motzkin path with $k$
up-steps, such that the path is positive after step $j$ for every corner
$(j,\eta_j)$ of $\eta$.

The same formula holds for every connected skew shape $\lambda/\mu$ with $n$
rows, where the board consists of the boxes $(j,c)$ with
$\mu_j<c<\lambda_{j+1}$, both for the lattice path matroid and for the
\hyperref[rookMatroids]{rook matroid} of $\lambda/\mu$.  The two matroids have
the same valuative invariants \cite{BoninMier2026}.

\begin{proof*}[Sketch]
The base polytope of a lattice path matroid has a subdivision, cut out by
integral prefix sums, whose maximal cells are snake matroid polytopes
\cite{BenedettiKnauerValencia2023x}, and $Z$ is a valuative invariant
\cite{FerroniSchroter2024}.  Every interior face of the subdivision is a
direct sum of connected ribbon matroids, which have Narayana $Z$-polynomials by
\hyperref[conjRibbonMatroidZ]{the ribbon theorem} below.  For a Ferrers shape,
the maximal cells correspond to the partitions $\eta\subseteq B$, and hence
to placements.  The faces of a fixed cell are obtained by fixing the
prefix sums at a subset of $m$ of its corner rows.  Such a face contributes
$(-1)^m$ times a product of the Narayana polynomials $Z_{U_{s,s+1}}$ of its
segments, which in gamma coordinates counts the Motzkin paths returning to zero
after each fixed row.  Inclusion--exclusion leaves
exactly the Motzkin paths that are positive at every corner row.
\end{proof*}


\subsection[schubertMatroidZRibbons]{Ribbons and outerplanar graphs}

A ribbon contains no $2\times2$ square, so the board is empty and the formula
above reduces to the Narayana polynomial.  This case is proved directly, and
is used in the proof above.

\begin{theorem}[Raj, 2026, \cite{Raj2026Outerplanarity}]
\label{conjRibbonMatroidZ}\label{ribbonMatroidZNarayana}
Every connected ribbon lattice path matroid $M$ of rank $r$ has
$Z_M(t)=N_{r+1}(t)$, independently of the composition of its row lengths.  In
particular, $Z_M(t)$ has $r$ simple negative zeros.
\end{theorem}

This was conjectured on the
\hyperref[realrootedOpenProblems]{open problems page}, and was known for the
circuits \cite{ProudfootXuYoung2018} and the fan matroids
\cite{LuXieYang2018x}.

\begin{proof*}[Sketch]
A connected ribbon matroid is graphic, represented by a path with a hub vertex
joined to both ends and to some interior vertices
\cite{KnauerMartinezSandovalAlfonsin2018}.  After removing parallel spokes,
it is an $(r+1)$-cycle with noncrossing chords at the hub.

Suppose a simple graph $G$ is the union of two graphs that meet exactly in an
edge $e=uv$ and its ends, and each side stays connected when $e$ is deleted.
The deletion formula of \name{Tom Braden} and \name{Artem Vysogorets}
\cite[Thm.~2.8]{BradenVysogorets2020} writes $Z_G-Z_{G\setminus e}$ as a sum of
multiples of $\tau(M/(F\cup e))$.  Here $\tau$ of a rank-$s$ matroid is the
coefficient of $t^{(s-1)/2}$ in its Kazhdan--Lusztig polynomial, and zero for
even $s$.  After contracting $e$ the two sides meet in a vertex, so
$M/(F\cup e)$ is a direct sum of two matroids of positive ranks $s_1,s_2$.
Since $P$ is multiplicative and $\deg P<\operatorname{rk}/2$ in positive rank,
the product has degree at most $(s_1+s_2-2)/2$, and $\tau$ vanishes.  Hence
$Z_G=Z_{G\setminus e}$.  Removing the chords one at a time leaves the cycle,
whose graphic matroid is $U_{r,r+1}$.
\end{proof*}

The same argument removes every diagonal of a dissected polygon, and the
converse holds.

\begin{theorem}[Raj, 2026, \cite{Raj2026Outerplanarity}]
Let $G$ be a simple $2$-connected graph on $n\geq3$ vertices.  Then $[t]Z_G(t)$
equals the number of bonds of $G$, and is at least $\binom{n}{2}$.  Moreover,
the following are equivalent: $G$ is outerplanar; $Z_G(t)=N_n(t)$;
$[t]Z_G(t)=\binom{n}{2}$.
\end{theorem}


\subsection[schubertMatroidZExample]{An example}

Take $\lambda=(3,3,2,1)$, so that $B$ has rows of lengths $2,1,0$.  Group the
placements by their set of occupied rows $J$.  The middle column is the
total weight of the placements with row set $J$, and the last column counts, by
up-steps, the Motzkin paths of length four that are positive after every step
in $J$.

\begin{tabular}{lll}
\toprule
$J$ & Rook weight & Motzkin paths \\
\midrule
$\emptyset$ & $1$ & $1+6y+2y^2$ \\
$\{1\}$ & $x_1+x_2$ & $3y+2y^2$ \\
$\{2\}$ & $x_1$ & $4y+y^2$ \\
$\{1,2\}$ & $x_1x_2$ & $2y+y^2$ \\
\bottomrule
\end{tabular}

At unit weights,
\[
  \Gamma_{M_\lambda}(y)=1+18y+8y^2,\qquad
  Z_{M_\lambda}(t)=1+22t+50t^2+22t^3+t^4.
\]
This is also the $Z$-polynomial of the staircase $S_4$ in the
\hyperref[conjNestedMatroidZ]{nested-matroid conjecture}, as the staircase
$(4,3,2,1)$ has the same board.


\subsection[schubertMatroidZWords]{A word formula}

The following formula needs only one family of words: the parity of a letter
gives the power of $t$, and its size decides whether a rook may be placed.

\begin{proposition}[Alexandersson, 2026]\label{schubertMatroidZWordFormula}
Let $\mathcal{A}_n$ be the set of words $a_1\dotsm a_n$ of nonnegative
integers with $a_i\leq a_{i-1}+1$ for all $i$, where $a_0=0$.  Then
\[
  Z_{M_\lambda}(t)=\sum_{a\in\mathcal{A}_n}
  t^{\#\{i\,:\,a_i\text{ odd}\}}\,
  W_B\bigl(\{j\in[n-1]:a_{j+1}\geq2\}\bigr).
\]
\end{proposition}

\begin{proof}
Start from a Motzkin path of length $n$ in the theorem above with some
horizontal steps marked.  Replace $U$ by $UU$, $D$ by $DD$, an unmarked $H$ by
$DU$ and a marked $H$ by $UD$, then add $U$ at the start and $D$ at the end.
This gives a \hyperref[qtCatalan]{Dyck path} of semilength $n+1$ whose height
at time $2j+1$ is $2h_j+1$.  Every Dyck path arises once, since it has
positive odd height at these times and the pairs of steps in between decode
uniquely.

Let $a_i$ be the height of the Dyck path just before its $(i+1)$st up-step,
for $0\leq i\leq n$.  Between two consecutive up-steps there are
$a_{i-1}+1-a_i$ down-steps, so the words $a_1\dotsm a_n$ are exactly the
elements of $\mathcal{A}_n$.  The $(i+1)$st up-step is at position $2i+1-a_i$.
It has even position precisely when $a_i$ is odd, and the even-position
up-steps come from the up-steps and the marked horizontal steps of the Motzkin
path.  Finally, $h_j>0$ means that the Dyck path has height at least three at
time $2j+1$, so that its $(j+2)$nd up-step is at position at most $2j+1$; this
says $a_{j+1}\geq2$.
\end{proof}

For example, for $n=2$ the board is empty, and the words $00,01,10,11,12$ give
$Z=1+3t+t^2$.

\begin{example}
For $n=4$, the Motzkin path $UHD\bar{H}$, with its last step marked, doubles to
the Dyck path $UUUDUDDUDD$.  The heights before its up-steps give the word
$a=1221$.  Its two odd letters come from the up-step and the marked step, and
the letters $a_2=a_3=2$ make rows $1$ and $2$ available, the rows where the
Motzkin path is positive.  The word contributes $t^2\,W_B(\{1,2\})$.
\begin{figure}
\svgimg[width=0.55\textwidth]{./svg-images/schubertZ-motzkin-dyck.svg}{The marked Motzkin path UHD with a marked last level step, and the Dyck path obtained by doubling it. Up-steps in even positions are highlighted, and each up-step is labelled by the height before it.}
\end{figure}
\end{example}  In terms of the Dyck path, $Z_{M_\lambda}$ is a sum over
\hyperref[standardYoungTableau]{standard Young tableaux} $T$ of shape
$(n+1,n+1)$, where the first row lists the up-step positions:
$T$ contributes $t$ to the power of the number of even entries in its first
row, and row $j$ of the board is available if $T_{1,j+2}\leq 2j+1$.

The words in $\mathcal{A}_n$ are the area sequences $(0,a_1,\dotsc,a_n)$ of
Dyck paths of semilength $n+1$, and the number of odd letters is the order of
vanishing at $q=-1$ of the product $\prod_i[a_i+1]_q$, which counts acyclic
orientations of a unit interval graph; see
\hyperref[dyckPathZPolynomials]{polynomials indexed by Dyck paths}.

The same words also give the gamma polynomial.  Write $a_i=2b_i+\epsilon_i$
with $\epsilon_i\in\{0,1\}$.  If $b_i\leq b_{i-1}$, both bits are possible, but
if $b_i>b_{i-1}$ then $b_i=b_{i-1}+1$ and $(\epsilon_{i-1},\epsilon_i)=(1,0)$
are forced.  Since $a_{j+1}\geq2$ precisely when $b_{j+1}>0$, the bits do not
affect the available rows, and summing over them gives the following.

\begin{corollary}[Alexandersson, 2026]
Let $\mathcal{B}_n$ be the set of words $b_1\dotsm b_n$ of nonnegative integers
with $b_1=0$, in which every rise $b_i>b_{i-1}$ has $b_i=b_{i-1}+1$ and no two
rises are consecutive.  Then
\[
  \Gamma_{M_\lambda}(y)=\sum_{b\in\mathcal{B}_n}
  y^{\#\text{rises of } b}\,
  W_B\bigl(\{j\in[n-1]:b_{j+1}>0\}\bigr).
\]
\end{corollary}

For example, $a=1234$ gives $b=0112$.  Its two rises fix all four bits, every
row is available, and the word contributes $t^2W_B(\{1,2,3\})$ to $Z$ and
$y^2W_B(\{1,2,3\})$ to $\Gamma$.


\subsection[schubertMatroidZPermutations]{Permutations, matchings and bases}

For a \hyperref[patternAvoidanceClasses]{$132$-avoiding} permutation
$\pi\in\symS_{n+1}(132)$, let $E(\pi)$ be the set of $j\in[n-1]$ such that the
values $1,\dotsc,n+1-j$ do not occupy a single interval of positions.  Then
\[
  Z_{M_\lambda}(t)=\sum_{\pi\in\symS_{n+1}(132)} t^{\des(\pi)}\,W_B(E(\pi)).
\]
To see this, append a final leaf to the Motzkin path and read $U,H,D$ in
preorder as vertices with two, one and zero children.  A unary vertex has its
child on the left for an unmarked step and on the right for a marked one.
Label the vertices $n+1,n,\dotsc,1$ in preorder and read them in inorder.  This
gives each $132$-avoiding permutation once, the descents come from the
vertices with a right child, and after $j$ vertices in preorder the $1+h_j$
unvisited subtrees are the intervals occupied by $1,\dotsc,n+1-j$.

\begin{example}
Let $n=3$ and $\lambda=(3,3,2)$, so that the board has rows of lengths $2$ and
$1$.  The fourteen permutations in $\symS_4(132)$ correspond to the fourteen
words in $\mathcal{A}_3$ and to the marked Motzkin paths of length three, where
$\bar{H}$ is a marked level step.  The last column is the number of
placements in the available rows.

\begin{tabular}{llllll}
\toprule
$a$ & $P$ & $\pi$ & $\des(\pi)$ & $E(\pi)$ & $W_B(E(\pi))$ \\
\midrule
$000$ & $HHH$ & $1234$ & $0$ & $\emptyset$ & $1$ \\
$001$ & $HH\bar{H}$ & $2134$ & $1$ & $\emptyset$ & $1$ \\
$010$ & $H\bar{H}H$ & $3124$ & $1$ & $\emptyset$ & $1$ \\
$100$ & $\bar{H}HH$ & $4123$ & $1$ & $\emptyset$ & $1$ \\
$012$ & $HUD$ & $2314$ & $1$ & $\{2\}$ & $2$ \\
$120$ & $UDH$ & $3412$ & $1$ & $\{1\}$ & $3$ \\
$122$ & $UHD$ & $2341$ & $1$ & $\{1,2\}$ & $5$ \\
$011$ & $H\bar{H}\bar{H}$ & $3214$ & $2$ & $\emptyset$ & $1$ \\
$101$ & $\bar{H}H\bar{H}$ & $4213$ & $2$ & $\emptyset$ & $1$ \\
$110$ & $\bar{H}\bar{H}H$ & $4312$ & $2$ & $\emptyset$ & $1$ \\
$112$ & $\bar{H}UD$ & $4231$ & $2$ & $\{2\}$ & $2$ \\
$121$ & $UD\bar{H}$ & $3421$ & $2$ & $\{1\}$ & $3$ \\
$123$ & $U\bar{H}D$ & $3241$ & $2$ & $\{1,2\}$ & $5$ \\
$111$ & $\bar{H}\bar{H}\bar{H}$ & $4321$ & $3$ & $\emptyset$ & $1$ \\
\bottomrule
\end{tabular}

Several rows give the same contribution $t^{\des(\pi)}W_B(E(\pi))$; adding
them up gives
\[
  Z_{M_\lambda}(t)=1+13t+13t^2+t^3=(1+t)(1+12t+t^2).
\]
\end{example}

Summing over signs gives a second formula.  Order signed integers as usual
and count descents between adjacent entries; then
\[
  Z_{M_\lambda}(t)=2^{-(n+1)}\sum_{\pi\in\symS_{n+1}(132)} W_B(E(\pi))
  \sum_{\epsilon\in\{\pm1\}^{n+1}} t^{\des(\epsilon\pi)}.
\]
A word with $p$ interior peaks contributes $2^{2p+1}t^p(1+t)^{n-2p}$ to the inner
sum, so equivalently
\[
  \Gamma_{M_\lambda}(y)=2^{-n}\sum_{\pi\in\symS_{n+1}(132)}W_B(E(\pi))\,
  (4y)^{\operatorname{pk}(\pi)},
\]
and $Z_{M_\lambda}$ is real-rooted if and only if this weighted peak polynomial
of $132$-avoiding permutations is.

Matching each up-step of a Motzkin path with its down-step gives a noncrossing
partial matching on $[n]$, and $h_j>0$ means that $j$ lies under an arc, that
is, $p\leq j<q$ for some arc $(p,q)$.  Thus $\gamma_k$ counts noncrossing
partial matchings with $k$ arcs, together with a placement whose rooks lie in
rows covered by arcs.

There is also a formula that starts from an actual basis rather than from the
board.  Read a basis $A$ of $M_\lambda$ in increasing order, and call
$b\in A$ \hyperref[internallyActive]{internally active} if no smaller element
$e\notin A$ makes $(A\setminus\{b\})\cup\{e\}$ a basis.  Then $\gamma_k$ counts pairs of
a basis $A$ and a Motzkin path with $k$ down-steps, in which an internally
active element is read as a step $U$ or $H$, and any other element as $U$, $H$
or $D$, but only from positive height.

Further descriptions of $Z_{M_\lambda}(t)$ in the same spirit sum
$t^{s}\,W_B(E)$ over the following objects, for the stated statistic $s$ and
row set $E$:
\begin{itemize}
\item lists $u_1<\dotsb<u_k$ and $v_1<\dotsb<v_k$ in $[n]$ with $u_i\leq v_i$,
  with $s=k$, and $j\in E$ if $u_i\leq j<v_i$ for some $i$;
\item plane rooted trees on $n+2$ vertices $v_0,\dotsc,v_{n+1}$, numbered in
  preorder, with $s$ the number of internal vertices minus one, and $j\in E$
  if $v_{n+1-j}$ is not on any of the paths that start at a child of the root
  and repeatedly take the last child;
\item permutations $\pi\in\symS_{n+1}(312)$, with $s=\des(\pi)$, and $j\in E$
  if position $j+1$ is the middle position of an occurrence of the pattern
  $231$.
\end{itemize}


\subsection[schubertMatroidZRowRecursion]{Adding a row and real-rootedness}

The word formula builds $Z$ one row at a time.  Appending a letter $q$ to a
word with last letter $p$ is allowed when $p\geq q-1$, contributes $t$ if $q$
is odd, and, if $q\geq2$, allows a rook in the new row of the board in any
column to the left of the last rook.

Let $\lambda$ have $m+1$ positive parts, let $Z_m$ be the $Z$-polynomial of
$(\lambda_1,\dotsc,\lambda_m)$, and let $Z_{m+1}$ be that of $\lambda$.  Let
$Q_{m+1}(t)$ be $t^{-1}$ times the sum in the word formula for $\lambda$ over
the words with last letter at least two; such a word contains an odd letter.
Every word for $(\lambda_1,\dotsc,\lambda_m)$ extends by the letters $0$ and
$1$ without making the new row available, so
\[
  Z_{m+1}(t)=(1+t)Z_m(t)+t\,Q_{m+1}(t).
\]

\begin{question}[Alexandersson, 2026]\label{schubertMatroidZRowInterlacing}
For every partition $\lambda$ with $m+1$ positive parts, does $Q_{m+1}(t)$
\hyperref[interlacingPolynomials]{interlace} $Z_m(t)$?
\end{question}

At a zero $r$ of $Z_m$, the recursion gives $Z_{m+1}(r)=rQ_{m+1}(r)$, so a
positive answer implies the
\hyperref[conjNestedMatroidZ]{nested-matroid conjecture} by induction on the
number of rows.  For $\lambda=(3,3,2,1)$ one has $Z_3=(1+t)(1+12t+t^2)$ and
$Q_4=8(1+3t+t^2)$.  At the zeros of $1+12t+t^2$, $Q_4=-72t>0$, while
$Q_4(-1)=-8$, so the zeros interlace.

Grouping the words is essential.  For $n=3$ and the empty board, the words
ending in $0$ and in $3$ give $1+3t+t^2$ and $t^2$, but $1+3t+3t^2$ is not
real-rooted, so the polynomials by last letter do not form an interlacing
family.

In the same unpublished work \cite{Alexandersson2026SchubertMatroidZ}, the
coefficient sequences of $Z_{M_\lambda}$ and $\Gamma_{M_\lambda}$ are shown to
be strictly log-concave for every Schubert matroid, and $Z_{M_\lambda}$ is
shown to be real-rooted for every Schubert matroid of rank at most $27$.
